\labsection{7}{Генетические алгоритмы с вещественным и перестановочным представлением особей. Задача коммивояжёра}

\labpart{Цель работы}
Изучить представление особей в форме вещественных чисел и перестановок, специализированные операторы
скрещивания и мутации для каждого представления, особенности решения задач непрерывной и комбинаторной
оптимизации генетическими алгоритмами; реализовать ГА, решающий задачу коммивояжёра.

\labpart{Теоретические сведения}
Подраздел~\ref{sec:real} и раздел~\ref{sec:combinatorial} теоретической части: вещественная
рекомбинация (\ref{eq:intermediate}), BLX-$\alpha$ (\ref{eq:blx}), гауссова и неравномерная мутации;
задача коммивояжёра (\ref{eq:tsp}), представления тура, операторы PMX, OX, CX, мутации перестановок,
локальный поиск 2-opt (\ref{eq:2opt}), меметические алгоритмы.

\labpart{Постановка задачи}
\emph{Часть А.} Найти минимум функции двух переменных своего варианта (таблица~\ref{tab:l6-var})
генетическим алгоритмом с представлением особей в форме вещественных чисел и сравнить результат с
бинарным ГА из ЛР~№6 при одинаковом бюджете вычислений функции.

\emph{Часть Б.} Дано множество городов, заданных координатами $(x_i; y_i)$ на плоскости; стоимость
перехода между городами равна евклидову расстоянию. Построить генетическим алгоритмом тур обхода всех
городов минимальной длины, используя представление путей, оператор кроссинговера и тестовый пример из
таблицы~\ref{tab:l7-var}. Сравнить найденный тур с оптимальным из файла \code{*.opt.tour} и
графически проиллюстрировать поиск.

\begin{table}[!htb]
\caption{Варианты заданий к лабораторной работе №7 (часть Б)}
\label{tab:l7-var}
\centering\small
\begin{tabular}{clccc}
\toprule
№ & Оператор кроссинговера & Файл TSPLIB & Городов & Оптимальный тур \\
\midrule
1 & частично соответствующий (PMX) & \code{berlin52.tsp} & 52  & 7542 \\
2 & упорядоченный (OX)             & \code{eil51.tsp}    & 51  & 426 \\
3 & циклический (CX)               & \code{eil76.tsp}    & 76  & 538 \\
4 & частично соответствующий (PMX) & \code{pr76.tsp}     & 76  & 108\,159 \\
5 & упорядоченный (OX)             & \code{st70.tsp}     & 70  & 675 \\
6 & циклический (CX)               & \code{lin105.tsp}   & 105 & 14\,379 \\
7 & частично соответствующий (PMX) & \code{kroA100.tsp}  & 100 & 21\,282 \\
\bottomrule
\end{tabular}
\end{table}

Файлы TSPLIB~\cite{reinelt1991} текстовые: заголовок (\code{NAME}, \code{TYPE : TSP}, \code{DIMENSION}~---
число городов, \code{EDGE\_WEIGHT\_TYPE : EUC\_2D}~--- евклидово расстояние, округлённое до целого),
раздел \code{NODE\_COORD\_SECTION} со строками <<номер $x$ $y$>> и строка \code{EOF}. В файле
\code{*.opt.tour} в разделе \code{TOUR\_SECTION} перечислены города оптимального тура, список
завершается числом $-1$; пересчёт его длины программой проверяет чтение данных и функцию расстояния.

\labpart{Требования к программе}
Исходные данные: размер популяции, вероятности кроссинговера и мутации, максимальное число поколений; в
части Б~--- файл \code{*.tsp} со списком городов. В части А отображаются популяция на линиях уровня
функции и графики лучшего и среднего значений функции. В части Б отображаются лучший тур текущей
популяции в виде графа на плоскости и в виде списка городов, графики лучшей и средней длины тура по
поколениям и длина лучшего тура. Отладку проводить на файлах TSPLIB, сравнивая решения с оптимальными
турами.

\labpart{Порядок выполнения работы}
\begin{steps}
\item Часть А: реализовать на NumPy вещественный ГА (турнирная селекция, BLX-$\alpha$ или
промежуточная рекомбинация, гауссова или неравномерная мутация с убывающим шагом, элитизм), выполнить
основной запуск с визуализацией; эталон найти \code{scipy.optimize.differential\_evolution}.
\item Исследовать влияние $\alpha$, шага и вероятности мутации, размера популяции (по 20 запусков при
сокращённом бюджете поколений, чтобы различия были заметны).
\item Сравнить вещественный и бинарный ГА из ЛР~№6 при одинаковых $N$ и числе поколений.
\item Часть Б: реализовать чтение файлов \code{*.tsp} и \code{*.opt.tour}, матрицу расстояний с
округлением по правилу \code{EUC\_2D} и длину тура (подраздел~\ref{sec:py-tsp}); проверить, что длина опубликованного оптимального тура
совпадает с указанной в таблице~\ref{tab:l7-var}.
\item Реализовать ГА с представлением путей: оператор кроссинговера варианта, мутация обменом или
инверсией, турнирная селекция, элитарное сокращение без дубликатов. Проверить оператор кроссинговера на
примере из подраздела~\ref{sec:combinatorial}.4 (для PMX и OX~--- также сравнением с DEAP).
\item Выполнить серию запусков и сравнить лучшую найденную длину с оптимальной.
\item Дополнительно: добавить локальный поиск 2-opt к каждому потомку (меметический ГА) и сравнить с
мультистартом 2-opt без ГА при одинаковом времени работы.
\end{steps}

\labpart{Пример выполнения}
\emph{Часть А.} Вариант~3: минимум 7-й функции Шаффера на $[-100; 100]^2$, эталон $f(0; 0) = 0$.
Вещественный ГА: $N = 60$, турнир $k = 3$, BLX-$0,5$ с $P_c = 0,9$, гауссова мутация ($P_m = 0,2$ на ген,
$\sigma_0 = 0,1$, шаг убывает линейно), две элиты. За 150 поколений (9060 вычислений функции) найдено
$f = 2,2 \cdot 10^{-9}$, порог $f < 10^{-3}$ пройден на 54-м поколении (рисунок~\ref{fig:l7-f7}).

Сравнение с бинарным ГА из ЛР~№6 при равном бюджете $60 \times 150$ (20 запусков): бинарный ГА
упирается в предел своей сетки (медиана и лучшее $f = 0,020$, худшее $0,18$), вещественный доходит до
медианы $3,5 \cdot 10^{-9}$ (лучшее $4,4 \cdot 10^{-12}$, худшее $0,015$). Влияние параметров при
сокращённом бюджете в 50 поколений приведено в таблице~\ref{tab:l7-params}.

\begin{figure}[!htb]
\centering
\includegraphics[width=0.92\textwidth]{\figdir/lab7a_population.png}
\caption{Популяция вещественного ГА на линиях уровня функции Шаффера F7 (масштаб осей уменьшается)}
\label{fig:l7-f7}
\end{figure}


\begin{table}[!htb]
\caption{Вещественный ГА на функции F7, 50 поколений: доля запусков с $f < 10^{-2}$ и медиана $f$
(20 запусков на значение)}
\label{tab:l7-params}
\centering\small
\begin{tabular}{lccc|lccc}
\toprule
 & Знач. & Успех, \% & Медиана $f$ & & Знач. & Успех, \% & Медиана $f$ \\
\midrule
$\alpha$ (BLX) & 0    & 20  & $4,3 \cdot 10^{-2}$ & $P_m$ & 0    & 95  & $1,3 \cdot 10^{-5}$ \\
               & 0,25 & 70  & $5,4 \cdot 10^{-3}$ &       & 0,05 & 95  & $5,1 \cdot 10^{-5}$ \\
               & 0,5  & 95  & $3,3 \cdot 10^{-3}$ &       & 0,2  & 95  & $3,3 \cdot 10^{-3}$ \\
               & 1    & 15  & $1,6 \cdot 10^{-2}$ &       & 0,5  & 0   & $1,9 \cdot 10^{-1}$ \\
\midrule
$\sigma_0$     & 0    & 100 & $1,5 \cdot 10^{-4}$ & $N$   & 20   & 5   & $1,8 \cdot 10^{-1}$ \\
               & 0,02 & 95  & $3,0 \cdot 10^{-3}$ &       & 40   & 55  & $6,6 \cdot 10^{-3}$ \\
               & 0,1  & 95  & $3,3 \cdot 10^{-3}$ &       & 60   & 95  & $3,3 \cdot 10^{-3}$ \\
               & 0,3  & 85  & $4,7 \cdot 10^{-3}$ &       & 120  & 100 & $1,8 \cdot 10^{-3}$ \\
\bottomrule
\end{tabular}
\end{table}

Главный оператор на этой функции~--- BLX-$\alpha$: при $\alpha = 0$ популяция сжимается раньше, чем доходит
до центра колец, при $\alpha = 1$ разброс потомков слишком велик для доводки. F7 радиально симметрична,
BLX-$0,5$ сам обеспечивает исследование, поэтому при малом бюджете крупный шаг мутации лишь мешает
доводке, а при $P_m = 0,5$ не успешен ни один запуск. Популяции из 20 особей недостаточно.

\emph{Часть Б.} Вариант~3: набор eil76 (76 городов, оптимальный тур длиной 538, средняя длина
случайного тура~--- 2520), циклический кроссинговер. Результаты приведены в таблице~\ref{tab:l7-tsp} и на
рисунке~\ref{fig:l7-tours}.

\begin{table}[!htb]
\caption{Задача коммивояжёра eil76, оптимум 538 (ГА~--- $150 \times 500$, меметический~--- $60 \times 60$)}
\label{tab:l7-tsp}
\centering\small
\setlength{\tabcolsep}{4pt}
\begin{tabular}{lccccc}
\toprule
Метод & Запусков & Лучший & Медиана & Худший & Отклонение, \% \\
\midrule
ГА: CX + обмен & 10 & 772 & 818 & 847 & $+52,0$ \\
ГА: CX + инверсия & 10 & 876 & 965 & 1071 & $+79,4$ \\
Меметический ГА & 5 & 538 & 538 & 538 & 0 \\
Мультистарт 2-opt & 5 & 540 & 550 & 553 & $+2,2$ \\
\bottomrule
\end{tabular}
\end{table}

\begin{figure}[!htb]
\centering
\includegraphics[width=\textwidth]{\labfig{7}{B1_tours.png}}
\caption{Лучшие найденные маршруты и опубликованный оптимальный тур eil76}
\label{fig:l7-tours}
\end{figure}

ГА с CX за 75\,000 оценок даёт туры на 43--57\,\% длиннее оптимального: CX сохраняет позиции городов,
а для маршрута важна их смежность, и рёбра хороших родителей почти не наследуются. Меметический ГА нашёл тур длиной 538 во всех пяти запусках
примерно за 5~с (маршрут отличается от опубликованного: у eil76 несколько оптимальных туров), а
мультистарт 2-opt за то же время~--- лишь 540--553.

\labpart{Контрольные вопросы}
\begin{questions}
\item Каковы преимущества вещественного кодирования? Опишите дискретную, промежуточную и линейную
рекомбинацию, BLX-$\alpha$.
\item Опишите равномерную, гауссову и неравномерную мутации вещественных особей. Зачем уменьшать шаг
мутации?
\item Сформулируйте задачи о рюкзаке, о покрытии и о коммивояжёре и их постановки для решения ГА.
\item Приведите способы представления тура и их особенности. Почему для задачи коммивояжёра не подходит
обычный двоичный кроссинговер?
\item Опишите представление соседства и операторы обмена рёбер, обмена подтуров и эвристического
скрещивания.
\item Опишите порядковое представление тура. Какой оператор кроссинговера для него применим?
\item Опишите операторы PMX, OX и CX для представления путей; выполните их для родителей
$(1\ 2\ 3\ 4\ 5\ 6\ 7\ 8)$ и $(3\ 7\ 5\ 1\ 6\ 8\ 2\ 4)$.
\item Какие мутации применяются к перестановкам? Опишите алгоритм 2-opt. Что такое меметический алгоритм и почему он эффективнее чистого ГА?
\end{questions}
